\chapter{Method}
\label{ch:method}

\section{Method Overview}
\label{sec:method_overview}

This chapter defines a method with four stages for predicting building attributes from (\acf{SVI}) and evaluating their effects on archetype assignment and registered energy class prediction.
The unit of analysis and linkage is one building recorded in the \acf{BAG}, the Dutch national register of addresses and buildings.
Its identifier is stored as \texttt{pand\_id}.
Building attributes refer to building type, construction year, and floor count.
Reference building attributes are obtained from EP-Online, \acs{BAG}, and reconstructed 3DBAG data, while predicted building attributes are produced by the vision configurations.
The Dutch national building typology developed within the \acf{TABULA} project is referred to as TABULA-NL \parencite{loga2016,tabulawebtool}.

The stages follow the information flow shown in Figure~\ref{fig:pipeline_overview}.
Stage 1 constructs the experimental dataset from reference records and selected street view images.
Stage 2 predicts building type, construction year, and floor count from the images.
Stage 3 combines building type and construction year to assign a TABULA-NL archetype cell and its representative U-values.
Stage 4 predicts the registered energy class either from building attributes and assigned U-values or directly from imagery.
This sequence separates model prediction from deterministic archetype assignment and uses the \acs{BAG} building as the unit of comparison throughout.

\begin{figure}[htbp]
	\centering
	\includegraphics[width=\linewidth]{figures/fig/fig_3_1_workflow}
	\caption[Four-stage method.]{Four-stage method for dataset construction, building attribute prediction, archetype assignment, and energy class prediction.}
	\label{fig:pipeline_overview}
\end{figure}

\subsection{Data Provenance}
\label{sec:three_layer_structure}

Throughout the four stages, the integrated dataset records the provenance of three field groups: source data, model outputs, and archetype properties.
Source data comprise administrative \acs{BAG} and EP-Online records and reconstructed 3DBAG geometry \parencite{kadasterbag,rvoeponline,threedbagdocs,peters2022,dukai2021}.
The administrative fields come from registry records. The 3DBAG fields describe geometry reconstructed from \acs{BAG} and national elevation data.
Model outputs comprise the predicted building attributes and the corresponding vision configuration.
Archetype properties comprise the construction period, TABULA-NL cell, and representative U-values assigned from either reference or predicted building type and construction year.
Every building assigned to the same archetype cell receives the same representative properties \parencite{loga2016,tabulawebtool}.
The registered EP-Online energy class remains the prediction target and is not an output of the archetype lookup.

\section{Stage 1: Dataset Construction}
\label{sec:dataset_construction}

Stage 1 constructs the experimental dataset required for all subsequent operations.
The reference branch produces one record per \acs{BAG} building that passes the reference filters, while the image branch produces one record per selected image.
The branches are linked by the \acs{BAG} building identifier \texttt{pand\_id}, so one building reference record may correspond to several street view image records.
Figure~\ref{fig:data_linkage_units} defines this record structure, and Table~\ref{tab:data_sources} identifies the contribution of each source.

\subsection{Data Sources}
\label{sec:data_sources}

The \acs{BAG} provides the building identifier, construction year, status, use, and two-dimensional footprint \parencite{kadasterbag}.
EP-Online is the Dutch national database of registered energy performance records and provides residential building category, registered energy class, and certificate registration date \parencite{rvoeponline}.
3DBAG is an open three-dimensional building dataset reconstructed from \acs{BAG} and national elevation data.
This study uses its geometric attributes rather than treating them as surveyed measurements \parencite{threedbagdocs,peters2022,dukai2021}.
Table~\ref{tab:3dbag_fields} defines the source fields and derived 3DBAG variables used in the analysis.

The image branch uses building footprints from \acf{OSM} and Overture Maps to locate candidate buildings, and Mapillary provides the associated panoramic imagery \parencite{openstreetmap,overturemaps,mapillary}.
TABULA-NL provides the archetype definitions and representative thermal properties used after attribute prediction \parencite{loga2016,tabulawebtool}.
All spatial operations use the Dutch national coordinate reference system, EPSG:28992.

\begin{figure}[htbp]
	\centering
	\includegraphics[width=\linewidth]{figures/fig/fig_3_2_record_structure_uml}
	\caption[Simplified Stage 1 record structure.]{Stage 1 record structure. Each \texttt{BuildingReferenceRecord} represents one \acs{BAG} building that passes the reference filters and is linked through \texttt{pand\_id} to one to eight \texttt{StreetViewImageRecord} entries. Each image record represents one selected facade crop.}
	\label{fig:data_linkage_units}
\end{figure}

\begin{table}[htbp]
	\centering
	\caption{Data sources, principal data elements, and roles in this study.}
	\label{tab:data_sources}
	\small
	\begin{tblr}{
		width = \linewidth,
		colspec = {Q[l,m,2.6cm] X[l] X[l]},
		colsep = 4pt,
		row{1} = {font=\bfseries},
		rowsep = 2pt,
	}
		\toprule
		Source & Principal data elements & Role in this study \\
		\midrule
		\acs{BAG} & Building identifier, construction year, two-dimensional footprint, status, and use & Linkage key, reference construction year, and residential building filter \\
		3DBAG & Reconstructed floor count, height, volume, and envelope areas at \acf{LoD} 2.2 & Reference floor count and geometric inputs for the transmission heat transfer calculation \\
		EP-Online & Residential building category, registered energy class, and registration date & Reference building type, prediction target, and certificate selection \\
		\acs{OSM} and Overture Maps & Building footprints and source identifiers & Candidate building locations and footprint linkage \\
		Mapillary & Panoramic street view images and image metadata & Visual input, image selection, and provenance \\
		TABULA-NL & Six construction periods and four residential size classes & Archetype cell assignment and representative as-built U-values \\
		\bottomrule
	\end{tblr}
\end{table}

\subsection{Building Reference Record Construction}
\label{sec:reference_dataset_construction}

The reference branch constructs one reference record for each BAG building. The normalised 16-digit identifier \texttt{pand\_id} is used to join BAG construction year and footprint data to EP-Online records and 3DBAG geometry. Within EP-Online, only records for existing residential buildings with a registered energy class are retained. When several certificate records are linked to one \texttt{pand\_id}, the record with the most recent registration date is selected. The registration date indicates whether the selected record was registered during the period governed by NTA 8800, the Dutch method for determining building energy performance \parencite{NTA8800}.

The reference floor count is taken from the reconstructed 3DBAG floor attribute. If this attribute is unavailable, floor count is estimated by dividing the median reconstructed building height by 3.0\,m and rounding the result to the nearest integer. Additional eligibility criteria for construction year, floor count, and TABULA-NL size class are listed in Table~\ref{tab:data_criteria}. Together, these steps produce one building reference record per \texttt{pand\_id} containing the reference building attributes, reconstructed geometry, and selected registered energy class.

\subsection{Street View Image Acquisition, Filtering, and Building Linkage}
\label{sec:svi_acquisition_matching}

The image branch adapts the acquisition and image processing workflow introduced by \textcite{liang2025} for OpenFACADES.
The workflow uses \acs{OSM} as its primary footprint source and Overture Maps to supplement missing footprints.
Mapillary supplies the 360\textdegree{} panoramas, and visibility analysis identifies candidate views of each building.
This study adds a \acs{BAG} residential filter to the image acquisition workflow.

Mapillary was selected because its licence permits an open and reproducible image acquisition workflow \parencite{mapillary}.
Its contributor-based coverage limits the evaluated buildings and does not represent a random sample of the Dutch residential building stock.

Candidate footprints and \acs{BAG} polygons are projected to EPSG:28992.
The centroid of each candidate footprint is intersected with \acs{BAG} polygons, and only footprints linked to a residential \acs{BAG} building proceed to image acquisition.
Each accepted crop is linked to the corresponding \acs{BAG} building through \texttt{pand\_id}, which joins the image record to the eligible reference set.

For each accepted footprint, Mapillary panorama metadata are queried within the study area shows in Figure~\ref{fig:svi_image_transformation}.
Candidate views are screened by camera distance, unobstructed visibility to the target building, and \acf{AoV}.
Views meeting the thresholds in Table~\ref{tab:data_criteria} are ranked, limited at footprint level, and downloaded at 2,048 pixels.
Grounding DINO, an open-set object detector that localises objects from textual prompts, identifies the target building using the prompt \texttt{building} \parencite{liu2024groundingdino}.

\begin{figure}[htbp]
	\centering
	\includegraphics[width=\linewidth]{figures/fig/fig_3_3_image_transformation}
	\caption[Image transformation from a 360\textdegree{} panorama to facade crops.]{Image transformation from a selected 360\textdegree{} panorama to facade crops for one building in Rotterdam. Imagery \textcopyright{} Mapillary contributors under the Creative Commons Attribution ShareAlike 4.0 licence (CC BY-SA 4.0).}
	\label{fig:svi_image_transformation}
\end{figure}

Crops are screened by building pixel share, wall pixel share, image sharpness, file size, and scene type.
Scene type is assessed with the ZenSVI implementation of Places365, a scene recognition model trained on 365 place categories.
Only crops classified as outdoor are accepted \parencite{zhou2018places,ito_zensvi_2025}.
Crops without a valid \acs{BAG} identifier are excluded.
The remaining crops are ranked by decreasing \ac{AoV} and increasing camera distance and are limited at building level according to Table~\ref{tab:data_criteria}.

\subsection{Final Dataset Integration and Eligibility}
\label{sec:svi_acquisition_filtering}
\label{sec:reference_data_integration}

Final integration joins the eligible reference set to the \acs{SVI} manifest by \texttt{pand\_id}.
The linked data product used in the experiments is termed the experimental dataset.
It is restricted to buildings with at least one image meeting the Stage 1 selection rules, processable 3DBAG geometry, and reference building attributes that permit TABULA-NL assignment.
All images linked to one building inherit the same reference building attributes and registered energy class.
Table~\ref{tab:data_criteria} provides the authoritative summary of the selection thresholds and outputs used throughout Stage 1.

\begin{table}[htbp]
	\centering
	\caption{Stage 1 selection rules and outputs.}
	\label{tab:data_criteria}
	\small
	\begin{tblr}{
		width = \linewidth,
		colspec = {Q[l,m,1.8cm] Q[l,m,3.1cm] X[l]},
		colsep = 4pt,
		row{1} = {font=\bfseries},
		rowsep = 3pt,
	}
		\toprule
		Branch & Process & Selection rule and output \\
		\midrule
		Reference & Building eligibility & Existing residential \acs{BAG} building with a registered energy class, construction year from 1800 to 2025, mapped residential category, and floor count from 1 to 30 \\
		Reference & Record linkage & \acs{BAG}, EP-Online, and 3DBAG records linked by normalised \texttt{pand\_id}, with one reference record selected per building \\
		Image & Footprint filter & \acf{OSM} or Overture Maps footprint centroid contained within a residential \acs{BAG} polygon \\
		Image & Panorama matching & Camera distance no greater than 30\,m and unobstructed visibility to the target building \\
		Image & View selection & $10^\circ < \acs{AoV} < 130^\circ$, with at most three panoramas per source footprint ranked by decreasing \acs{AoV} \\
		Image & Crop quality & Building pixel share at least 15\%, wall pixel share no more than 30\%, Laplacian variance above 30, file size at least 20\,kB, and classification as outdoor by the Places365 scene classifier \\
		Image & Building image limit & Valid \acs{BAG} assignment and at most eight images per building, ranked by decreasing \acs{AoV} and increasing camera distance \\
		Integration & Experimental dataset & Reference record and at least one image meeting the Stage 1 selection rules, processable 3DBAG geometry, and valid TABULA-NL input attributes \\
		\bottomrule
	\end{tblr}
\end{table}

\subsection{Data Partitioning and Cross-Validation}
\label{sec:data_partitioning}

The experimental dataset is partitioned before model fitting. This procedure separates model development from final evaluation in two steps.

First, the scikit-learn \texttt{StratifiedGroupKFold} splitter generates five shuffled subsets with random seed 42 while attempting to balance the distribution of a composite stratification label \parencite{pedregosa2011}.
One subset is designated as the fixed holdout set, and the remaining four subsets form the development set.
The development set is used for model fitting, validation, and model selection.
The fixed holdout set is the predefined partition reserved for final evaluation and is not used for model fitting, hyperparameter selection, or checkpoint selection.
The composite stratification label combines study area, original EP-Online residential category, registered energy class, and TABULA-NL construction period.
The split is generated from one record per building, using \texttt{pand\_id} as the grouping variable, and all images linked to that building inherit its assignment.
Images of the same building therefore cannot occur in different subsets.
This grouping prevents repeated images of one building from crossing subsets, but it does not remove spatial dependence between neighbouring buildings, which can make random-split performance optimistic \parencite{garbasevschi2021}.

Second, the same splitter divides the development set into five folds using the same stratification variables, grouping key, and random seed. Each trainable model is fitted five times, with four folds used for training and the remaining fold used for validation in each run. This procedure provides validation results for model and checkpoint selection without using the holdout set.
The same predefined partitions are used for building attribute prediction, archetype assignment, and energy class prediction.

\section{Stage 2: Building Attribute Prediction}
\label{sec:vision_inference_strategies}

Stage 2 predicts building type, construction year, and floor count from the street view images linked to each \acs{BAG} building.
Building type is predicted using four classes, whose labels supply the corresponding TABULA-NL residential size class during Stage 3.
Construction year and floor count are treated as continuous prediction targets during supervised training.
The Stage 2 output records the vision configuration and predicted building attributes for each building.
ResNet-50 and DINOv2 produce one set of three attributes for every building.
InternVL3-2B produces a building-level prediction only when at least one image response can be parsed and all three values fall within the permitted ranges.
Reporting the three predicted attributes separately permits their evaluation before building type and construction year enter the deterministic TABULA-NL lookup.

\subsection{Vision Configurations}
\label{sec:vision_model_configurations}

Three vision configurations represent different degrees of adaptation to the development data.
ResNet-50 provides a convolutional neural network whose residual architecture supports full supervised fine tuning \parencite{he2016}.
DINOv2 provides self-supervised visual representations through a \acf{ViT} backbone that can remain frozen while new prediction layers are trained \parencite{oquab2023}.
InternVL3-2B provides a pretrained vision language model that can return building attribute predictions from images and a text prompt without updating its parameters \parencite{zhu2025}.

Table~\ref{tab:vision_configs} defines the three complete configurations.
They differ in architecture, pretraining, parameter updates, number of images, and aggregation procedure.
The comparison therefore evaluates the complete configurations and does not isolate any one of these design choices.
The two supervised configurations process at most eight images per building, whereas InternVL3-2B processes at most three because it performs a separate generative inference for each image.
In the two supervised configurations, one \acf{MLP} transforms the aggregated image representation before it enters the three attribute-specific output heads.

\begin{table}[htbp]
	\centering
	\caption{Vision configurations used for building attribute prediction.}
	\label{tab:vision_configs}
	\footnotesize
	\begin{tblr}{
		width = \linewidth,
		colspec = {Q[l,m,2.7cm] X[l] X[l] Q[c,m,1.3cm] X[l]},
		colsep = 3pt,
		row{1} = {font=\bfseries},
		rowsep = 2pt,
	}
		\toprule
		Configuration & Parameters updated on the development set & Development labels used & Image limit & Building-level aggregation \\
		\midrule
		ResNet-50 full fine tuning & Backbone, \acs{MLP}, and three attribute heads & All training fold labels & 8 & Mean of image feature vectors before prediction \\
		DINOv2 \acs{ViT}-B/14 with frozen backbone & \acs{MLP} and three attribute heads & All training fold labels & 8 & Mean of image feature vectors before prediction \\
		InternVL3-2B without parameter updating & None & Prompt design using a fixed subset of 200 development buildings & 3 & Type mode, year median, and rounded floor count median \\
		\bottomrule
	\end{tblr}
\end{table}

\subsection{Supervised Attribute Prediction}
\label{sec:supervised_attribute_prediction}

ResNet-50 and DINOv2 use the same building-level prediction structure.
Each available image is transformed to $224 \times 224$ pixels and converted by the relevant backbone into a feature vector.
Feature vectors from at most eight images are averaged to produce one representation per building.
The representation passes through the \acs{MLP} with a hidden dimension of 256 and then through separate heads for building type, construction year, and floor count.

The joint training objective combines cross entropy (CE) loss for building type with mean squared error (MSE) loss for construction year and floor count:

\begin{equation}
	\label{eq:attr_loss}
	\mathcal{L}_{\mathrm{attr}}
	= \mathcal{L}_{\mathrm{CE,type}}
	+ \mathcal{L}_{\mathrm{MSE,year}}
	+ \mathcal{L}_{\mathrm{MSE,floor}} .
\end{equation}

The cross entropy term uses weights inversely proportional to class frequency in the training fold.
Construction year and floor count are standardised using the mean and standard deviation of that fold before the two mean squared error terms are calculated.
The three losses receive equal weight.
ResNet-50 updates its backbone, \ac{MLP}, and output heads, while DINOv2 keeps the backbone frozen and updates only the \ac{MLP} and output heads.
The optimiser, learning rate, augmentation, and early stopping settings are reported in Table~\ref{tab:vision_hyperparams}.

Both supervised vision configurations are fitted separately in each development fold.
The validation metric is the macro F1 score for building type, calculated as the unweighted mean of the F1 scores for the four classes.
This metric controls early stopping and selects one checkpoint within each fold.
The checkpoint used on the holdout set is selected only from validation performance in the development data.
No holdout labels are used to fit parameters, choose hyperparameters, or select checkpoints.

\subsection{Building Attribute Prediction with InternVL3-2B}
\label{sec:vlm_attribute_inference}

InternVL3-2B uses the pretrained \texttt{OpenGVLab/\allowbreak InternVL3-2B} checkpoint without parameter updates on the study data \parencite{zhu2025}.
The prompt instructs the model to return one JavaScript Object Notation object per image with building type, construction year, and floor count.
Greedy decoding is used so repeated inference with the same input and environment follows a fixed decoding rule.
The final prompts, valid ranges, and parsing rules are fixed before the holdout set is processed and are reported in Appendix~\ref{ch:appendix_d}.

InternVL3-2B processes the three highest-ranked images available for a building independently.
A response is valid only when all three required values can be parsed and fall within the permitted ranges.
Building type is aggregated by the mode.
Construction year is aggregated by the median, and floor count by the rounded median.
No building-level prediction is created when all image responses are invalid.

\section{Stage 3: Archetype Assignment and Thermal Parameters}
\label{sec:tabula_archetype_matching}

Stage 3 uses a fixed lookup to assign a TABULA-NL archetype cell based on building type and construction year. Building type determines the residential size class, and construction year determines the construction period. Their joint assignment selects the representative as-built U-values for the wall, roof, floor, and windows. The same TABULA-NL lookup is applied to the reference and predicted values of building type and construction year. Therefore, any difference between the resulting archetype cells is caused by differences in these input values, not by different lookup rules.

\subsection{TABULA-NL Lookup Structure}
\label{sec:tabula_nl_matrix}

The TABULA-NL lookup used in this study contains six construction periods and four residential size classes, forming 24 archetype cells, as shown in Figure~\ref{fig:tabula_matrix}.
The period codes and thermal properties are generated programmatically from the \acs{TABULA} values workbook and compared with the \acs{TABULA} WebTool \parencite{tabulawebtool}.
The generation rules and the identified difference between the workbook and WebTool are documented in Section~\ref{sec:appendix_a_tabula}.

\begin{figure}[htbp]
	\centering
	\includegraphics[width=0.88\linewidth]{figures/fig/fig_3_6_tabula_matrix}
	\caption[Operational TABULA-NL lookup.]{TABULA-NL lookup used in this study with construction periods and residential size classes. Each cell shows a representative dwelling and its archetype identifier. Adapted from the \acs{TABULA} WebTool \parencite{tabulawebtool}.}
	\label{fig:tabula_matrix}
\end{figure}

The four size classes are Single-Family House (SFH), Terraced House (TH), Multi-Family House (MFH), and Apartment Block (AB).
No additional rule based on dwelling count or floor count is used to determine these classes.
Table~\ref{tab:type_mapping} defines how the selected EP-Online categories are converted to the four TABULA-NL classes.

\subsection{Mapping Attributes to Archetypes}
\label{sec:attribute_archetype_mapping}

Let $s_i$ denote the TABULA-NL size class of building $i$, and let $p(y_i)$ denote the construction period obtained by applying the TABULA-NL year boundaries to construction year $y_i$.
The archetype cell is

\begin{equation}
	\label{eq:archetype_cell}
	c_i = \bigl(s_i,\, p(y_i)\bigr).
\end{equation}

For assignment from reference building attributes, the \acs{BAG} construction year supplies $y_i$ and the EP-Online building type is converted to $s_i$ using Table~\ref{tab:type_mapping}.
For assignment from predicted building attributes, the predicted construction year supplies $y_i$ and the predicted building type supplies $s_i$ directly.
Both input sets use the same construction period boundaries and lookup table.
Floor count is included in energy class prediction but does not enter Equation~\ref{eq:archetype_cell}.

\begin{table}[htbp]
	\centering
	\caption{Mapping from EP-Online residential categories to TABULA-NL size classes.}
	\label{tab:type_mapping}
	\small
	\begin{tblr}{
		width = \linewidth,
		colspec = {X[l] X[l] Q[l,m,2.8cm]},
		colsep = 4pt,
		row{1} = {font=\bfseries},
		rowsep = 2pt,
	}
		\toprule
		EP-Online category & English description & TABULA-NL size class \\
		\midrule
		Vrijstaande woning & Detached house & SFH \\
		Twee-onder-één-kap & Semi-detached house & SFH \\
		Rijwoning tussen & Mid-terrace house & TH \\
		Rijwoning hoek & End-terrace house & TH \\
		Woongebouw met niet-zelfstandige woonruimte & Residential building with non-self-contained accommodation & MFH \\
		Appartement & Apartment & AB \\
		\bottomrule
	\end{tblr}
\end{table}

\subsection{Representative U-Value Assignment and As-Built Assumption}
\label{sec:uvalue_lookup}

After cell assignment, the lookup returns one U-value in $\mathrm{W/(m^2\,K)}$ for each of four envelope elements: wall, roof, floor, and windows.
Table~\ref{tab:tabula_uvalues} reports the values used by the implementation.
Single-Family House, Multi-Family House, and Apartment Block share the same values within a construction period, while Terraced House differs in period NL.01.

\begin{table}[htbp]
	\centering
	\caption{TABULA-NL as-built envelope U-values ($\mathrm{W/(m^2\,K)}$).}
	\label{tab:tabula_uvalues}
	\small
	\begin{tblr}{
		width = \linewidth,
		colspec = {Q[l,m,1.5cm] Q[l,m,3.0cm] Q[l,m,2.4cm] Q[r] Q[r] Q[r] Q[r]},
		colsep = 4pt,
		row{1} = {font=\bfseries},
		rowsep = 2pt,
	}
		\toprule
		Period & Construction year range & TABULA-NL size class & $U_{wall}$ & $U_{roof}$ & $U_{floor}$ & $U_{window}$ \\
		\midrule
		NL.01 & $\leq$1964 & SFH/MFH/AB & 1.6129 & 1.5385 & 1.7241 & 5.2 \\
		NL.01 & $\leq$1964 & TH & 2.2222 & 2.0833 & 2.4390 & 5.2 \\
		NL.02 & 1965 to 1974 & All & 1.4493 & 0.8929 & 2.3256 & 5.2 \\
		NL.03 & 1975 to 1991 & All & 0.6410 & 0.6410 & 1.2821 & 5.2 \\
		NL.04 & 1992 to 2005 & All & 0.3584 & 0.3584 & 0.3584 & 2.9 \\
		NL.05 & 2006 to 2014 & All & 0.2660 & 0.2347 & 0.2660 & 1.8 \\
		NL.06 & $\geq$2015 & All & 0.2101 & 0.1597 & 0.2660 & 1.8 \\
		\bottomrule
	\end{tblr}
\end{table}

Renovation status is not available consistently for individual buildings, so every assignment uses the unrenovated, as-built variant of its TABULA-NL cell.
Buildings with the same size class and construction period consequently receive the same four U-values \parencite{loga2016}. The lookup does not account for renovations or measured envelope properties of individual buildings.
They are used as inputs to the energy class classifier and are combined with reconstructed 3DBAG envelope areas for the transmission heat transfer comparison.
The U-values describe thermal transmittance of individual envelope elements.
When combined with reconstructed element areas and divided by estimated floor area, they yield the floor-area-normalised transmission heat transfer coefficient $h_{tr}$ used in this study \parencite{loga2013calculation}.

\section{Stage 4: Energy Class Prediction}
\label{sec:downstream_model}

Stage 4 implements two procedures for predicting the registered energy class.
The first uses building attributes, assigned TABULA-NL U-values, and study area as structured inputs.
The second maps street view images directly to the energy class without producing intermediate attributes or an archetype cell.
Both procedures use the same predefined building partitions and target definitions.
Comparisons among energy class conditions use the evaluation set defined in Section~\ref{sec:dataset_statistics}.

EP-Online records the registered energy class of each dwelling.
The residential scale runs from G, the lowest-performing class, to A++++, the highest-performing class, with additional plus signs distinguishing performance beyond class A \parencite{rvoenergielabel}.
The primary prediction target groups classes A to C and D to G.
For labels determined under NTA 8800, this division corresponds to 250\,kWh/(m$^2\cdot$year), the upper limit of class C \parencite{rvoenergielabel2024}.

\subsection{Energy Class Prediction from Building Attributes}
\label{sec:attribute_based_energy_class_prediction}

The structured prediction procedure uses \acf{LightGBM}, a tree-based gradient boosting framework configured for classification \parencite{ke2017}.
Its eight inputs are building type, construction year, floor count, four envelope U-values, and study area.
Building type and study area are categorical, while the remaining inputs are numerical.
Other certificate fields are excluded because they would not be available for a building whose energy class is to be estimated.

For each target, one \acs{LightGBM} classifier is fitted on the complete development set using reference attributes and their assigned U-values.
The fitted classifier is applied first to equivalent reference inputs in the holdout set and then to predicted attributes with U-values assigned from the predicted type and year.
Study area remains unchanged.
This design holds the classifier constant and limits the input difference to the three building attributes and the U-values affected by type and year.

Both target definitions use 400 boosting iterations, a learning rate of 0.05, 31 leaves, and no class weighting.
No separate hyperparameter search is performed for the two holdout input sets.
The complete \acs{LightGBM} configuration is reported in Table~\ref{tab:lightgbm_hyperparams}.

\subsection{Direct Prediction of Energy Class from Images}
\label{sec:direct_energy_class_prediction}

The direct image procedure predicts the registered energy class without first estimating building attributes.
It follows earlier work that mapped building imagery directly to registered energy performance classes \parencite{mayer2023uk,mayer2023dk}.

The direct ResNet-50 configuration uses the same image processing, feature aggregation, and parameter update settings as its attribute prediction configuration, but replaces the three attribute heads with one energy class head.
The direct DINOv2 configuration reuses the pretrained DINOv2 ViT-B/14 backbone from the attribute prediction configuration without updating its parameters. Feature vectors from the images of each building are averaged and cached, and a separate energy class prediction network is trained on these building representations.
Separate configurations are fitted for the two targets using cross entropy loss and class weights from the training fold.
Validation macro F1 controls early stopping and checkpoint selection.
Section~\ref{sec:appendix_b_direct_energy} reports the corresponding optimisation settings.

InternVL3-2B uses a separate fixed prompt for each target without parameter updates.
The building class is selected by majority vote over valid responses from the three highest-ranked images.
Ties are resolved in favour of the more energy efficient class or class group.
Stage 4 returns one energy class prediction per building, target, and vision configuration.
